\section{The remaining mathematical construction}
\label{sec:remaining}

The analytic conclusion and the arithmetic construction have different
statuses. The preceding sections prove that suitable estimates for
corrected prime sums imply the Riemann hypothesis and Dirichlet GRH.
They also give many algebraic, local analytic and lattice comparisons.
The missing step is to supply their remaining quantitative inputs for
one compatible family of integral maps and geometric objects. This
section states that obligation by mathematical dependency.

\subsection{The final analytic acceptance condition}

Fix a primitive target character $\chi$. At each accuracy stage choose
a finite positive character combination $H_j$ retaining $\chi$ or
$\bar\chi$, and keep that combination fixed during its sequence of
arbitrarily late intervals. By
Theorem~\ref{thm:scale-averaged-criterion}, it is enough that
$\theta_j\downarrow1/2$ and that the weighted positive-part integral
\eqref{eq:scale-averaged-target} have subpower growth along such
intervals. The interval exponent and local averaging width are fixed
at each stage; they may change between stages.

The arithmetic comparison must hold almost everywhere on the same
expanded intervals used by that local averaging. Signed contributions
may be combined and locally averaged before the positive part is taken.
An uncontrolled signed integral is insufficient. Pointwise estimates
throughout whole intervals remain a stronger sufficient supplier.
For the numerical construction it suffices that its total exponent
loss tends to zero, as in Corollary~\ref{cor:variable-loss-construction};
no uniform bound across accuracy stages is required. Bounds inside a
fixed stage must still accommodate the actual growing ranks and fields.

\subsection{The geometric degree and its integral corrections}

The required component lattices carry prescribed Petersson metrics
and periods. Corollary~\ref{cs-mean-volume-intersection} reduces
the mean-volume input to one normalized total arithmetic intersection
bound and the prescribed metric and section-counting comparison.
Its finite-prime correction is elementary. The auxiliary volume limit
fixes each datum before taking section powers to infinity; no uniform
finite convergence threshold is needed in that step. The intersection
bound for the selected growing family remains a geometric obligation.

Lemma~\ref{geometric-volume-weights} sums the exact
volume comparison before discarding its component degrees.
Proposition~\ref{geometric-one-mean} uses the short integral sections
to remove the need for an independent negative-height moment.
Corollary~\ref{geometric-one-mean-scale} then gives a uniform degree
estimate from the genuine mean-volume input and the stated
Hilbert--Samuel estimate. The flat twist is charged once, afterward.

Alternatively, Proposition~\ref{geometric-signed-raw-degree} retains
the actual section ranks and signed remainder. Its final cost contains
no component heights and imposes no separate upper rank-share
hypothesis. That explicit cost still needs a bound for the actual
growing family. Neither its smallness nor the mean-volume theorem
follows from the existence of short sections alone.

These degrees must then pass through the actual covering and type
maps. Corollary~\ref{geometric-cover-improved} retains their rank ratio.
Corollary~\ref{cs-exceptional-budget} records the discarded degree
and the integral splitting index, including their signs. A small
retained character family does not make these indices bounded.
Every lattice, metric, period and projection used here must agree
with the arithmetic identity.

\subsection{Integral theta maps and centered global energy}

A theta lift is an automorphic integral obtained from the Weil
representation for a pair of groups. The proposed construction uses
chains between unitary groups over CM fields, which are totally
imaginary quadratic extensions of totally real fields. The chain
length and the auxiliary fields may grow with the scale.
The integral coefficient modules, pairings, levels and measures must
be specified for these actual chains. Their required nonvanishing
and sharp auxiliary-level divisibility remain part of the construction.
The local coefficient and period comparisons established earlier do
not alone supply those global assertions.

Proposition~\ref{cs-exact-energy} gives a smaller sufficient frame
interface: the diagonal energy $\Delta$ and centered global energy
$V$ of the actual normalized maps. Separate inequalities for every
pair of words and the large redundancy used by that supplier are not
intrinsic requirements of the slope theorem. The exact positive
normalizers and the identification with the same integral maps in a
common target basis must be retained. Restriction to the permitted
source subspaces is controlled by
Proposition~\ref{cs-retained-energy}; the target rank-loss denominator
remains separate.

The centered comparison of Proposition~\ref{cs-centered-slope} has
zero loss for a perfect full frame. It may be used with the optimum
normalization of Proposition~\ref{cs-optimal-normalization}, taking
the minimum with the earlier valid bounds. Its square-root dependence
on the target energy need not improve every numerical regime.
Exact proportional expansion would remove the source-spread and
minimum-slope assumptions by Corollary~\ref{cs-exact-expansion},
but that expansion has not been proved for the constructed family.
These comparisons concern the same quotient sources and final target;
they do not remove the intermediate geometric degree or the separate
counting, divisor and quotient errors.

The centered kernel integral in Section~\ref{cs-averaged-theta-energy}
is a concrete global quantity to estimate. Its total mass alone does
not control it. Local orbital identities, finite-word residue positivity
and the ideal-summation formulas are available ingredients, but a
uniform global passage to this energy is still required. The stronger
orbital conditions \eqref{th-four-estimates} and
\eqref{th-diag-input} remain possible suppliers. The divisibility
condition \eqref{cs-required-divisor} must concern the maps in this
same energy calculation. Corollary~\ref{cs-refined-combined-budget}
then keeps the correlation between energy errors, source slopes and
target heights. A joint estimate over maps and scale is sufficient
when it supplies the final weighted budget.

\subsection{Cohomology, selection and compatibility}

The classical Galois representations and their local and infinity
conditions are supplied by
Theorem~\ref{thm-classical-galois-zero-monodromy}, with common period
realizations as in Lemma~\ref{lem-galois-period-realizations}.
This replacement removes a separate supplier of those representations;
it does not remove independent theta identities or geometric estimates.
Rational adjoint vanishing does not by itself bound integral congruence
lengths.

The cohomological estimates must use the full decomposition-group
conditions at the coefficient prime and include torsion and primes
dividing the auxiliary degree. The coarse common-interval mechanism
uses Theorem~\ref{kt-fourteen}; the opposite-parity estimate is in
Theorem~\ref{ms-main}. Remaining complementary-parity automorphic
components require estimates in the allowed growing fields.
The spectral and multiplicity formulas must also yield nonempty
selected spaces with the asserted ranks and constants. Existing
formulas do not replace those quantitative conclusions.

Under its stated uniform quasi-Riemann-hypothesis input,
Theorem~\ref{qs-selection} supplies simultaneous quadratic selection
on the required finite family. Section~\ref{sec:selection-assembly}
fixes the compatible order of parameters. These inputs concern the
specified Dirichlet characters; they are not automatically estimates
for arbitrary Hecke characters of varying CM fields. All conductor,
rank, word and exceptional-space costs must be included in the same
final discrepancy.

\subsection{The ordinary Riemann hypothesis}

For the conductor-one target the analytic function is always $G_1$,
even when auxiliary characters and geometric data change with the
scale. The positive primitive-character kernel in
Corollary~\ref{rh-direct-budget} compares the larger weights before
averaging. The larger-argument quasi-Riemann-hypothesis input controls
the remaining samples, including the corrected principal character.
No symmetry of the weight distributions is required.

Corollary~\ref{rh-averaged-final} leaves one precise target: a weighted
positive-part integral of the locally averaged sum of the normalized
arithmetic discrepancy and its signed quadratic companion. The
transfer and centering remainders already have finite weighted tail
integrals. A separate companion estimate is one sufficient supplier,
but the combined signed sum may be used directly. The auxiliary
geometric ranks and integral structures still grow, so the same
geometric, theta-energy and compatibility obligations remain.

\begin{example}[A fixed zero-free region and symmetry do not suffice]
Let $\vartheta>1/2$, choose
$0<a<\min(\vartheta-1/2,1/2)$, and put
\[
 P(s)=\bigl((s-\tfrac12-a)^2+1\bigr)
       \bigl((s-\tfrac12+a)^2+1\bigr).
\]
Then $P(1-s)=P(s)$ and
$P(\bar s)=\overline{P(s)}$. It has no real zero and no zero
with $\Re s>\vartheta$, but has four zeros off the critical
line, at $1/2\pm a\pm i$. Every assertion follows by direct
factorization. This polynomial is not a Dirichlet $L$-function
and is not offered as a counterexample to GRH. It shows that
a fixed zero-free region, absence of real zeros, and the usual
reflection symmetries are logically insufficient by themselves.
The arithmetic features of the $L$-functions must supply the
remaining argument.
\end{example}

The missing compatible arithmetic estimate must be proved, or replaced
by an independently proved supplier of the averaged criterion. Until
then this manuscript is a collection of rigorous reductions and stated
conditional constructions, not a completed proof of RH or Dirichlet GRH.
